\documentclass[a4paper]{article}
% Español (México) con XeLaTeX: fontspec para fuentes Unicode, polyglossia
% para la división silábica y los nombres del documento en la variante mexicana.
\usepackage{fontspec}
\usepackage{polyglossia}
\setdefaultlanguage[variant=mexican]{spanish}
% Los nombres chinos van como «pinyin (original)»: xeCJK da a los caracteres
% Han su propia fuente; sin ella XeLaTeX los omite con un aviso.
% FandolSong no cubre algunos caracteres de nombres (U+73FA); WenQuanYi Zen Hei sí.
\usepackage[AutoFallBack=true]{xeCJK}
\setCJKmainfont{FandolSong-Regular.otf}
\setCJKfallbackfamilyfont{\CJKrmdefault}{WenQuanYi Zen Hei}
% polyglossia no traduce los nombres de listings.
\AtBeginDocument{\providecommand{\lstlistingname}{}\renewcommand{\lstlistingname}{Listado}%
  \providecommand{\lstlistlistingname}{}\renewcommand{\lstlistlistingname}{Índice de listados}}
\usepackage{amsmath, amssymb}
\usepackage{graphicx}
\usepackage[margin=2.5cm]{geometry}
\usepackage[most]{tcolorbox}
\usepackage{etoolbox}
\usepackage{listings}
\usepackage{hyperref}
\usepackage{booktabs}
\usepackage{subcaption}
\usepackage{float}

\newtcolorbox{knowledgebox}[1]{enhanced, colback=blue!5!white, colframe=blue!75!black, colbacktitle=blue!75!black, coltitle=white, fonttitle=\bfseries, title={#1}, attach boxed title to top left={yshift=-2mm, xshift=2mm}, boxrule=1pt, sharp corners}
\newtcolorbox{importantbox}[1]{enhanced, colback=yellow!10!white, colframe=yellow!80!black, colbacktitle=yellow!80!black, coltitle=black, fonttitle=\bfseries, title={#1}, sharp corners}
\newtcolorbox{warningbox}[1]{enhanced, colback=red!5!white, colframe=red!75!black, colbacktitle=red!75!black, coltitle=white, fonttitle=\bfseries, title={#1}, sharp corners}

\lstset{language=Python, basicstyle=\ttfamily\small, keywordstyle=\color{blue}, stringstyle=\color{red!60!black}, commentstyle=\color{green!60!black}, breaklines=true, frame=single, numbers=left, numberstyle=\tiny\color{gray}, captionpos=b, extendedchars=true}

\newcommand{\notetitle}{{[}LLM Agents F24{]} Measuring Agent Capabilities and Anthropic's RSP — Ben Mann}
\newcommand{\noteauthors}{Elaboradas a partir de materiales públicos del curso}
\newcommand{\notedate}{25 de noviembre de 2024}
\newcommand{\videochannel}{Berkeley RDI}
\newcommand{\videopublishdate}{25 de noviembre de 2024}
\newcommand{\videoduration}{aproximadamente 60 minutos}
\newcommand{\videourl}{https://www.youtube.com/live/6y2AnWol7oo}
\newcommand{\videocoverpath}{cover.jpg}
\newcommand{\repourl}{https://github.com/hqhq1025/ai-course-notes}


% Footer with repo info
\usepackage{fancyhdr}
\pagestyle{fancy}
\fancyhf{}
\fancyfoot[C]{\thepage}
\fancyfoot[R]{\footnotesize\color{gray}\href{\repourl}{AI Course Notes}}
\renewcommand{\headrulewidth}{0pt}
\renewcommand{\footrulewidth}{0pt}

\begin{document}
\begin{titlepage}\centering
{\Large Notas del curso\par}\vspace{1.2cm}
{\huge\bfseries \notetitle\par}\vspace{0.8cm}
{\large \noteauthors\par}\vspace{0.3cm}
{\large \notedate\par}\vspace{1.2cm}
\ifdefempty{\videocoverpath}{{\small Al generar las notas, indique la ruta local de la portada del video.\par}}{\includegraphics[width=0.82\textwidth,height=0.45\textheight,keepaspectratio]{\videocoverpath}\par}
\vfill
\begin{tcolorbox}[width=0.9\textwidth, colback=black!2!white, colframe=black!60, sharp corners]
\textbf{Autor o canal del video}: \videochannel\par
\textbf{Fecha de publicación}: \videopublishdate\par
\textbf{Duración del video}: \videoduration\par
\ifdefempty{\videourl}{}{\textbf{Enlace del video}: \href{\videourl}{\nolinkurl{\videourl}}\par}\par
\textbf{Página del proyecto}: \href{\repourl}{\nolinkurl{\repourl}}\par
\end{tcolorbox}
\end{titlepage}

\tableofcontents\newpage


\section{Introducción: Anthropic y la seguridad en la IA}

Ben Mann es uno de los cofundadores de Anthropic y fue uno de los primeros autores del artículo de GPT-3. Anthropic se fundó en 2021 con el propósito de construir sistemas de IA seguros. Esta conferencia se centra en dos preguntas fundamentales: ¿cómo \textbf{medir} la capacidad de un agente de IA? ¿Cómo formular una \textbf{política de escalamiento responsable} a partir de la evaluación de esa capacidad?

\begin{knowledgebox}{La línea de investigación de Anthropic}
La investigación de Anthropic abarca: la interpretabilidad de circuitos en transformers (entender los mecanismos internos del modelo), Constitutional AI (hacer que la IA se autocorrija con base en principios), el desarrollo y despliegue de la serie de modelos Claude, y la formulación de la Responsible Scaling Policy.
\end{knowledgebox}

\subsection{Resumen de la sección}

La idea central de Anthropic es: cuanto mayor sea la capacidad de la IA, más estrictas deben ser las medidas de seguridad. Medir la capacidad es la condición previa para formular las medidas de seguridad.

\section{Evaluación de las capacidades de los Agent}

\subsection{Por qué es necesario evaluar las capacidades de los Agent}

\begin{importantbox}{Doble propósito de la evaluación}
\begin{enumerate}
    \item \textbf{Conocer el avance}: cuantificar el nivel de capacidad del Agent en distintos tipos de tareas
    \item \textbf{Anticipar riesgos}: identificar las dimensiones de capacidad que se acercan a un umbral peligroso, para desplegar medidas de seguridad con anticipación
\end{enumerate}
\end{importantbox}

\subsection{Dimensiones de la evaluación}

\begin{itemize}
    \item \textbf{Autonomía}(Autonomy): ¿puede el Agent completar tareas de ciclo largo sin supervisión humana?
    \item \textbf{Uso de herramientas}(Tool Use): ¿con qué eficacia usa el Agent herramientas externas (ejecución de código, búsqueda, API)?
    \item \textbf{Capacidad de persuasión}(Persuasion): ¿puede el Agent influir en las decisiones humanas? Esto implica un riesgo potencial de manipulación.
    \item \textbf{Conocimiento CBRN}: ¿el Agent posee conocimiento especializado en los ámbitos químico, biológico, radiológico y nuclear?
    \item \textbf{Ciberseguridad}: ¿puede el Agent descubrir y explotar vulnerabilidades de software?
\end{itemize}

\subsection{Metodología de evaluación}

\begin{itemize}
    \item Diseñar tareas del mundo real (no puntos de referencia sintéticos) y observar el desempeño del Agent de principio a fin
    \item Comparación con una línea base humana: hacer que expertos del dominio realicen las mismas tareas y comparar el desempeño del Agent
    \item Prestar atención a las «capacidades de cola»: las capacidades que el modelo podría mostrar en situaciones extremas
\end{itemize}

\begin{warningbox}{Dificultad fundamental de la evaluación}
La evaluación es simétrica: si podemos evaluar una capacidad, es porque ya la entendemos. Pero lo más peligroso puede ser precisamente aquello que aún no hemos identificado (unknown unknowns).
\end{warningbox}

\subsection{Resumen de la sección}

La evaluación de las capacidades de los Agent debe abarcar múltiples dimensiones, como la autonomía, el uso de herramientas, la influencia social y el conocimiento especializado, y debe ir más allá de los puntos de referencia estándar para prestar atención al desempeño en el mundo real y a los riesgos de cola.

\section{Política de Escalamiento Responsable (RSP)}

\subsection{El marco central de la RSP}

\begin{importantbox}{Niveles ASL (AI Safety Level)}
Anthropic propuso un sistema de clasificación por niveles de seguridad en IA (AI Safety Level), similar a los niveles de bioseguridad (BSL):
\begin{itemize}
    \item \textbf{ASL-1}: sistemas de IA sin riesgo significativo
    \item \textbf{ASL-2}: nivel de los grandes modelos actuales, requiere medidas de seguridad estándar
    \item \textbf{ASL-3}: posee capacidad autónoma o conocimiento especializado significativos, requiere medidas de seguridad reforzadas
    \item \textbf{ASL-4+}: cercano a capacidades sobrehumanas, requiere medidas de seguridad extremas
\end{itemize}
Cada nivel corresponde a requisitos de seguridad distintos: condiciones de despliegue, intensidad de monitoreo, restricciones de uso.
\end{importantbox}

\subsection{El mecanismo de funcionamiento de la RSP}

\begin{enumerate}
    \item \textbf{Evaluación periódica}: cada vez que se entrena un modelo nuevo o se realiza una actualización significativa, se evalúa si su capacidad cruza un límite de ASL
    \item \textbf{Mecanismo de activación}: si la evaluación muestra que el modelo se acerca a la capacidad del siguiente nivel, se suspende el despliegue hasta que estén listas las medidas de seguridad correspondientes
    \item \textbf{Inversión en seguridad}: vincular el presupuesto de investigación y desarrollo en seguridad con el crecimiento de la capacidad del modelo
    \item \textbf{Transparencia}: publicar periódicamente informes de evaluación de capacidades
\end{enumerate}

\subsection{El significado de la RSP}

\begin{knowledgebox}{Por qué importa la RSP}
La contribución central de la RSP es vincular explícitamente la ``capacidad del modelo'' con las ``medidas de seguridad''; no se trata de ``publicar primero y remediar después'', sino de ``asegurar primero y avanzar después''. Esto le da a toda la industria una plantilla de marco de seguridad operativa.
\end{knowledgebox}

\subsection{Resumen de la sección}

Mediante la clasificación por niveles ASL y el mecanismo de activación, la RSP logra una ``escalada de seguridad impulsada por la capacidad''. Este es, hasta ahora, uno de los marcos de seguridad en IA más operativos.

\section{Fideicomiso de Beneficio a Largo Plazo}

Anthropic creó de manera innovadora el Fideicomiso de Beneficio a Largo Plazo (Long-term Benefit Trust): a medida que la IA se vuelva cada vez más poderosa, el control de la empresa se transferirá gradualmente a directores independientes sin interés económico, cuyo único enfoque sea la seguridad y el interés público. Esta es la primera estructura de gobernanza de este tipo en la industria de la IA.

\begin{importantbox}{Innovación en gobernanza}
En un contexto donde la IA podría llegar a ser una de las tecnologías más poderosas de la historia, asegurar que los incentivos de quienes ejercen el control estén alineados con el interés público es una garantía institucional para la seguridad a largo plazo.
\end{importantbox}

\section{Benchmarking Real-World Agents}

\subsection{Benchmark Pipeline}

Anthropic descompone la capacidad del Agent en varias métricas: autonomía, invocación de herramientas, profundidad de razonamiento, consistencia de preferencias. Cada ronda de evaluación encadena las siguientes etapas:

\begin{enumerate}
    \item \textbf{Design scenarios}: selección de tareas de ciclo largo (investigación, planeación, ejecución, reporte)
    \item \textbf{Establish baselines}: comparación con equipos humanos o modelos de generaciones anteriores
    \item \textbf{Instrument data}: recopilación de los registros de acciones del agent, las llamadas a la API y el trail de auditoría
    \item \textbf{Score metrics}: calificación automatizada + revisión humana
    \item \textbf{Feedback to training}: cada failure mode detectado se convierte, uno por uno, en un dataset/behavioral trace
\end{enumerate}

\begin{figure}[H]
\centering
\includegraphics[page=3,width=0.92\textwidth]{slides.pdf}
\caption{Agent Benchmark pipeline illustrated on slide 3.\protect\footnotemark}
\end{figure}
\footnotetext{Slides emphasize multi-step evaluation loops.}

\subsection{Case Study: Decision Support Task}

Una tarea típica consiste en que el Agent lea un documento interno de estrategia de 20 páginas, genere un resumen POAP (Plan, Options, Actions, Priority) y programe una reunión de follow-up.

\begin{itemize}
    \item El Agent necesita mantener la coherence entre páginas (anidamiento de párrafos, referencias)
    \item La salida debe atender a la vez los hechos, la cadena de razonamiento y las recomendaciones de acción
    \item Al final, un reviewer humano asigna la calificación según la template de `POAP`
\end{itemize}

\begin{knowledgebox}{Tailored judgment rubric}
Anthropic elabora un rubric para cada tipo de tarea: exactitud factual, integridad estructural, cumplimiento en la invocación de herramientas, si ofrece o no recomendaciones accionables. Solo se avanza a la actualización de ASL cuando todas las dimensiones son satisfactorias.
\end{knowledgebox}

\subsection{Resumen de la sección}

El benchmark pipeline de Anthropic encadena varios escenarios de tarea, la recopilación de datos, la calificación y la retroalimentación al entrenamiento en un ciclo cerrado; cada scenario cuenta con un rubric específico, con lo cual la evaluación de capacidades se convierte en una señal de gobernanza accionable.

\subsection{Evidence Matrix}

Para convertir la evaluación de capacidades en información para la toma de decisiones, Anthropic usa una \textbf{evidence matrix}: enumera, por etapa de la tarea, las métricas clave, las fuentes de datos, el desempeño actual, y da un límite de seguridad (green/orange/red).

\begin{table}[H]
\centering
\begin{tabular}{p{0.18\textwidth}p{0.24\textwidth}p{0.24\textwidth}p{0.24\textwidth}}
\toprule
Etapa & Métrica & Fuente de datos & Señal de riesgo actual \\
\midrule
Planning & clarity of plan, dependency awareness & prompt logs, human evaluation & plan misses key stakeholders \\
Tool invocation & tool success rate, error rate & API logs, retry counts & spike in `api\_error` events \\
Execution & prospect completion quality & manager reviews, outcome diff & 5\% of tasks reworked \\
Reporting & factual accuracy & summary diff + external documents & hallucinated reference detected \\
\bottomrule
\end{tabular}
\caption{Evidence matrix used for capability snapshots}
\end{table}

\begin{warningbox}{Evidence matrix refresh cadence}
Matrix rows are refreshed every week or after any significant dataset/model update. Without timely refresh, the safety gates rely on stale performance snapshots.
\end{warningbox}
\section{Monitoring, Telemetry, and Incident Response}

\subsection{Telemetry Stack}

Cada llamada a herramienta y cada salida de la API del Agent desplegado se escribe en la cadena de observabilidad, que fluye hacia tres tipos de sistemas:

\begin{itemize}
    \item \textbf{Alerting}: se establecen umbrales con base en anomaly detection, como despliegues automatizados repentinos o requests de alta frecuencia
    \item \textbf{Replay}: registra el comportamiento completo para su reproducción, incluye timeline, embeddings, entradas/salidas
    \item \textbf{Metrics}: métricas personalizadas (alignment drift, policy deviation, human override rate)
\end{itemize}

\begin{warningbox}{Observability isn't optional}
En cuanto el Agent se desvía del modo habitual, es necesario poder reproducir la cadena de llamadas completa y ver el system prompt y el context externo. Sin esta visualización, la investigación y la gobernanza son casi imposibles.
\end{warningbox}

\subsection{Incident Response Loop}

Cuando el monitoreo dispara una alerta, el proceso de respuesta de Anthropic incluye:

\begin{enumerate}
    \item Determinar rápidamente el nivel de riesgo (green/orange/red)
    \item Cambiar automáticamente a un modo de marioneta, en el que todas las salidas requieren aprobación humana
    \item Abrir un post-mortem, registrando el timeline y la root cause
    \item Actualizar el checklist y las instrucciones del prompt, para evitar la recurrence
\end{enumerate}

\subsection{Resumen de la sección}

El monitoreo del Agent necesita el respaldo de tres patas: telemetry, alert y replay; en cuanto se dispara una alerta, la organización debe responder con rapidez mediante un structured loop (determinación de nivel, containment, post-mortem, update) para evitar que la capacidad se deslice hacia la zona de riesgo.

\subsection{Data Pipelines and Instrumentation}

Telemetry data is ingested into a multi-tier pipeline: real-time streaming, batch aggregation, and backing datastore for audit. Each agent interaction emits structured events with the following fields:

\begin{itemize}
    \item `task\_id`, `agent\_variant`, `score`
    \item `tools\_used` list with duration/evidence per tool
    \item `alignment\_drift` delta vs baseline
    \item `human\_override` flag
\end{itemize}

\begin{importantbox}{Why instrumentation matters}
Instrumentation not only powers alerts but also enables retrospective capability audits: you can slice by agent version, time window, tool set, or base capability and replay the exact logs that triggered a red flag.
\end{importantbox}

The ingestion stack keeps both raw and derived layers. Raw events are immutable and tagged with `event\_hash` to ensure idempotency. Derived views pre-aggregate per ASL level and feed dashboards that teams monitor in daily standups.

Dashboards themselves include `alignment drift`, `tool success`, `red-team exposure`, and `human override latency`. When any metric crosses thresholds, the telemetry system automatically tickets the incident response crew.

\section{Deployment Guardrails and Governance}

\subsection{Pre-production Gates}

Los gates clave antes del despliegue incluyen:

\begin{itemize}
    \item \textbf{Red team test}: un equipo independiente busca posibles formas de uso indebido
    \item \textbf{Explainability check}: se confirma, mediante herramientas de interpretability, la evidence chain clave del model
    \item \textbf{Configuration hardening}: se bloquean el toolset, los ratelimits y los system prompts, y se registran en `CLAUDE.md`
\end{itemize}

\subsection{Governance Artifacts}

Anthropic mantiene un conjunto de artifact (policy doc, runbooks, SLAs) y establece que:

\begin{itemize}
    \item cada nueva entrega debe ir acompañada de un `Capability Report`
    \item el Stakeholder tiene derecho a request un `Transparency Brief`
    \item los trustees del `Long-term Benefit Trust` tienen veto
\end{itemize}

Artifacts are versioned in a governance repository and cross-referenced with the deployment pipeline. Every update must cite the specific evidence matrix row it addresses and include the slide page demonstrating the capability.

\subsection{Action Timeline}

\begin{figure}[H]
\centering
\includegraphics[page=7,width=0.92\textwidth]{slides.pdf}
\caption{Slide 7 outlines governance milestones and Trust handoff.\protect\footnotemark}
\end{figure}
\footnotetext{Trust handoff timeline with independent board oversight.}

\subsection{Resumen de la sección}

Los guardrails se componen de pre-production gates, artifact y el veto del Trust; solo si estas medidas se aplican en cada punto de incremento de capacidad se puede evitar que la seguridad de la IA se convierta en un remedio posterior.

\subsection{Action Timeline}

Deployment is treated as a sequence of checkpoints. A typical timeline:

\begin{tabular}{p{0.22\textwidth}p{0.64\textwidth}}
\toprule
Phase & Controls \\
\midrule
Pre-launch & red team, trust briefing, capability report, artifact sign-off, slides deck review \\
Soft launch & limited roll-out, parallel human monitor, data capture instrumentation, immediate rollback plan \\
Full launch & full monitoring, trustee oversight, public communications, community feedback loop \\
\bottomrule
\end{tabular}

\begin{enumerate}
    \item Pre-launch checklist includes red-team signoff, evidence matrix alignment, and trustee update.
    \item Soft launch extends agent access to a limited set of collaborators while instrumentation is verified.
    \item Full launch triggers marketing communication plus continuous monitoring with trustee review cadence.
\end{enumerate}

\begin{warningbox}{Action timelines must not shortcut}
Skipping any checkpoint (e.g., moving to soft launch before artifacts sign-off) is a common root cause for incidents. Trust briefer explicitly checks for missing evidence pages in `slides.pdf` before signing off.
\end{warningbox}

\section{Open Questions and Research Agenda}

\subsection{Measurement Gaps}

Anthropic planteó públicamente several open gaps:

\begin{itemize}
    \item How to quantify ``alignment drift'' in the wild?
    \item How to estimate human override fatigue when multiple Agent versions run concurrently?
    \item How to instrument slow-moving capabilities (e.g., strategic reasoning) before they surface in benchmarks?
\end{itemize}

\begin{knowledgebox}{Action Items}
Every new capability launch must include (1) measurement plan, (2) failure signature sheet, (3) human-in-the-loop trigger.
\end{knowledgebox}

\subsection{Research Directions}

Ben Mann highlights:

\begin{itemize}
    \item ``Capability contour'' research—mapping the Pareto frontier of autonomy vs. controllability
    \item Multimodal evaluation—ensuring Agent maintains alignment when tool outputs are visual/audio
    \item Policy orchestration—how to mix synchronous/asynchronous Agent families safely
\end{itemize}

\subsection{Resumen de la sección}

El mayor desafío de investigación en esta etapa es lograr que las tres líneas de measurement/safety/governance avancen de manera simultánea: se necesitan alignment metrics más detalladas, multi-agent orchestration studies, así como verificar la eficacia real del artifact de gobernanza.

\section{Metrics Glossary}

\subsection{Key Capability Metrics}

Anthropic tracks a catalog of capability metrics that feed into the evidence matrix. The main categories are:

\begin{itemize}
    \item \textbf{Integrity}: factual accuracy, hallucination rate, citation precision
    \item \textbf{Autonomy}: plan depth, tool use depth, self-corrections
    \item \textbf{Control}: human override ratio, rollback events, command adherence
    \item \textbf{Impact}: productive output, decision quality, human satisfaction index
\end{itemize}

Each metric is tied to a measurement method: `Integrity` uses reference corpora and automated fact-checkers, `Autonomy` uses action logs, `Control` uses human override logs, `Impact` uses structured surveys.

\subsection{Glossary Table}

\begin{tabular}{p{0.18\textwidth}p{0.22\textwidth}p{0.3\textwidth}p{0.22\textwidth}}
\toprule
Metric & Description & Data Source & Safety Threshold \\
\midrule
Fact Density & Share of claims verified by docs & embeddings + text diff & >82\% \\
Plan Completeness & Steps with dependency resolution & plan transcripts & no missing references \\
Tool Success & Tools used / tools invoked & telemetry logs & >95\% \\
Override Lag & Time between agent action and human override & workflow logs & >3 min to override \\
Drift Index & Deviation from baseline policy & alignment scores & <0.07 drift \\
Policy Noise & entropy in action selection & policy logs & <=0.3 \\
\bottomrule
\end{tabular}

\begin{warningbox}{Metric calibration}
Metrics must be recalibrated when models change architecture, since scores are relative to behavior and can drift even if underlying capability remains stable.
\end{warningbox}

\section{Operational Checklists}

\subsection{Incident Triage Checklist}

Anthropic operationalizes incidents using a multi-step checklist:

\begin{enumerate}
    \item Document the deviation and gather log artifacts
    \item Estimate likelihood/severity using ASL heuristics
    \item Engage second reviewer and determine containment actions
    \item Log the event in the incident tracker and assign owner
    \item Update capability notes/training data to reflect the failure
\end{enumerate}

\subsection{Deployment Playbooks}

Each deployment uses two playbooks—`Launch Playbook` and `Rollback Playbook`. The launch playbook enumerates gating evidence (matrix rows, telemetry dashboards, trustee review). The rollback playbook spells out evaluation metrics, check for repeated incursion, and communication plan.

\subsection{Resumen de la sección}

Metrics glossary + checklists keep tooling grounded: each number in the evidence matrix is tied to a data source, each incident follows a documented escalation path. These artifacts make evaluation defensible even under audit scrutiny.

\section{Slide Gallery}

The provided slides illustrate the evidence-heavy thinking behind these policies. Each slide screenshot is presented to anchor the textual summary above.

\begin{figure}[H]
\centering
\includegraphics[page=6,width=\textwidth,height=0.9\textheight]{slides.pdf}
\caption{Slide 6: capability contours and instrumentation.\protect\footnotemark}
\end{figure}
\clearpage
\begin{figure}[H]
\centering
\includegraphics[page=8,width=\textwidth,height=0.9\textheight]{slides.pdf}
\caption{Slide 8: telemetry architecture sketch.\protect\footnotemark}
\end{figure}
\clearpage
\begin{figure}[H]
\centering
\includegraphics[page=9,width=\textwidth,height=0.9\textheight]{slides.pdf}
\caption{Slide 9: incident response steps.\protect\footnotemark}
\end{figure}
\clearpage
\begin{figure}[H]
\centering
\includegraphics[page=11,width=\textwidth,height=0.9\textheight]{slides.pdf}
\caption{Slide 11: governance artifact checklist.\protect\footnotemark}
\end{figure}
\clearpage
\begin{figure}[H]
\centering
\includegraphics[page=13,width=\textwidth,height=0.9\textheight]{slides.pdf}
\caption{Slide 13: Trust transfer timeline.\protect\footnotemark}
\end{figure}
\clearpage
\begin{figure}[H]
\centering
\includegraphics[page=14,width=\textwidth,height=0.9\textheight]{slides.pdf}
\caption{Slide 14: measurement gaps \& research agenda.\protect\footnotemark}
\end{figure}
\clearpage
\begin{figure}[H]
\centering
\includegraphics[page=15,width=\textwidth,height=0.9\textheight]{slides.pdf}
\caption{Slide 15: operational checklist for deployments.\protect\footnotemark}
\end{figure}
\clearpage
\footnotetext{Slides exported from the provided deck; pages chosen to illustrate the written narrative.}

\section{Análisis operativos a fondo}

\subsection{Flujo de trabajo del informe del equipo rojo}

Después de cada mejora de capacidad, el equipo rojo elabora un informe narrativo. Incluye la cronología, la evidencia que detonó el hallazgo y los pasos de mitigación. La narrativa se estructura en (1) contexto previo al incidente, (2) vector de capacidad, (3) comportamiento observado, (4) modo de falla hipotetizado, (5) controles recomendados.

El consejo de gobernanza revisa el informe y lo incorpora al Capability Report. Cuando varios informes citan el mismo modo de falla, se crea un `Issue Cluster` y se asigna a un `Capability Owner` dedicado.

\subsection{Informes de cumplimiento y transparencia}

Anthropic publica periódicamente informes breves de transparencia que describen los cambios del modelo, los resultados de las pruebas y las preocupaciones pendientes. Cada informe se valida contra el calendario del consejo del Long-term Benefit Trust.

El flujo de trabajo de transparencia incluye tres aprobaciones: la del líder de investigación, la del líder de seguridad y la del fideicomisario. Cualquier desviación antes de obtener la aprobación desencadena una pausa y una explicación formal al consejo.

\subsection{Estrategia de registro de datos}

Todos los datos de instrumentación (métricas de capacidad, llamadas a herramientas, anulaciones) se transmiten de manera continua a un almacén multiusuario con registros `[tag, timestamp, event\_id, payload]`. El pipeline de ingesta etiqueta los eventos con `project`, `agent\_version`, `ASL\_level`.

Los registros sin procesar se conservan durante 90 días; los paneles agregados (desviación de la alineación, éxito de las herramientas) se conservan durante 365 días. Esto permite auditorías retrospectivas mucho después de un despliegue.

\subsection{Resumen de la sección}

Estos análisis a fondo muestran que el trabajo del equipo rojo, la transparencia y el registro de datos no son improvisados, sino que están entretejidos en rituales de gobernanza regulares. Cada ritual produce artefactos que alimentan la evaluación comparativa, el monitoreo y las barreras de seguridad del despliegue.
\section{Síntesis y ampliación}

\subsection{Tabla de síntesis}

\begin{table}[H]
\centering
\begin{tabular}{p{0.2\textwidth} p{0.26\textwidth} p{0.26\textwidth} p{0.16\textwidth}}
\toprule
Ámbito & Estrategia clave & Medios de implementación & Pregunta abierta \\
\midrule
Benchmarking & Evaluación basada en escenarios + rubric & Decision support case + POAP + slide rubric & ¿Cómo cuantificar el alignment drift? \\
Monitoring & Telemetry + replay + alerts & Metrics/alert logs + incident loop & ¿Cómo mantener la interpretabilidad durante la ejecución con varios agentes? \\
Governance & Pre-deploy gates + Trust oversight & red teams + capability reports + trustee veto & ¿Cómo integrar el governance artifact en el CI/CD? \\
Research & Capability contour + multi-agent orchestration & Long-term Benefit Trust + measurement plan & ¿Cómo capturar los tail-end risks before deployment? \\
\bottomrule
\end{tabular}
\caption{Tabla que resume las cuatro dimensiones de la Lecture 02: capacidad, seguridad, gobernanza e investigación}
\end{table}

\subsection{Lecturas adicionales}

\begin{itemize}
    \item Anthropic, ``Responsible Scaling Policy'' 2023
    \item Anthropic, ``Claude Model Cards and Capability Reports'' 2024
    \item Anthropic, ``Long-term Benefit Trust'' governance brief
    \item Ben Mann, talk materials from Berkeley RDI (slides included in repo)
    \item Shevlane et al., ``Model Evaluation for Extreme Risks,'' 2023
\end{itemize}

\subsection{Resumen de la sección}

La Lecture 02 enlaza la evaluación de capacidades, el monitoreo, los guardrails de despliegue y la gobernanza institucional en un ciclo cerrado: evaluación continua + monitoreo + pre-deploy gates + trustee oversight; la agenda de investigación, por su parte, lleva más allá el benchmark measurement, la policy orchestration y el governance artifact, sentando un marco de acción para la seguridad futura de los Agent.

\subsection{Slide Evidence}

The lecture slides underline the emphasis on evidence-backed governance. Slide 4 distills the capability matrix, while slide 5 captures the Trust handoff timeline.

\begin{figure}[H]
\centering
\includegraphics[page=4,width=0.92\textwidth]{slides.pdf}
\caption{Slide 4: Capability matrix and risk signals.\protect\footnotemark}
\end{figure}

\begin{figure}[H]
\centering
\includegraphics[page=5,width=0.92\textwidth]{slides.pdf}
\caption{Slide 5: Long-term Benefit Trust timeline.\protect\footnotemark}
\end{figure}
\footnotetext{Screenshots extracted from the provided slide deck.}

\end{document}
